% ============================================================================
% DIAGRAMES PER AL TFG CITYFIX
% ============================================================================
% Afegir al preàmbul del document principal:
%
% \usepackage{tikz}
% \usetikzlibrary{arrows.meta, positioning, shapes.geometric, fit,
%                  backgrounds, calc, decorations.pathreplacing}
% ============================================================================


% ============================================================================
% DIAGRAMA 1: MÀQUINA D'ESTATS XSTATE — CICLE DE VIDA D'INCIDÈNCIES
% ============================================================================

\begin{figure}[H]
\centering
\begin{tikzpicture}[
    node distance=2.8cm and 3.2cm,
    state/.style={
        rectangle, rounded corners=8pt, draw=black, thick,
        minimum width=2.6cm, minimum height=1.1cm,
        font=\sffamily\small, align=center
    },
    open/.style={state, fill=blue!15, draw=blue!60},
    assigned/.style={state, fill=yellow!20, draw=yellow!60!black},
    inprogress/.style={state, fill=orange!20, draw=orange!60},
    validated/.style={state, fill=green!15, draw=green!60!black},
    closed/.style={state, fill=gray!20, draw=gray!60},
    trans/.style={
        -Stealth, thick, font=\sffamily\scriptsize
    },
    guard/.style={
        font=\sffamily\tiny, text=gray!70!black, align=center
    },
    every edge quotes/.style={font=\sffamily\scriptsize, auto}
]

% Estats
\node[open]       (open)       {OPEN\\{\tiny Oberta}};
\node[assigned, right=of open]   (assigned)   {ASSIGNED\\{\tiny Assignada}};
\node[inprogress, right=of assigned] (inprog)     {IN\_PROGRESS\\{\tiny En procés}};
\node[validated, below=2.5cm of inprog]  (validated)  {VALIDATED\\{\tiny Validada}};
\node[closed, left=of validated]    (closed)     {CLOSED\\{\tiny Tancada}};

% Punt inicial
\node[circle, fill=black, inner sep=2pt, left=1.2cm of open] (init) {};
\draw[trans] (init) -- (open);

% Punt final
\node[circle, draw=black, thick, inner sep=2pt, below=0.8cm of closed] (final) {};
\node[circle, fill=black, inner sep=1.2pt] at (final) {};

% Transicions principals
\draw[trans, blue!70!black] (open) -- node[above, guard] {ASSIGN} node[below, guard] {\textit{isNotStudent}} (assigned);

\draw[trans, orange!70!black] (assigned) -- node[above, guard] {START} node[below, guard] {\textit{isNotStudent}} (inprog);

\draw[trans, green!60!black] (inprog) -- node[right, guard] {RESOLVE} node[left, guard, xshift=-2mm] {\textit{isTechnicalOrAdmin}} (validated);

\draw[trans, gray!70!black] (validated) -- node[above, guard] {CLOSE} node[below, guard] {\textit{isAdmin}} (closed);

\draw[trans] (closed) -- (final);

% Transicions inverses
\draw[trans, red!60!black, bend left=30] (validated) to node[right, guard] {REJECT} node[left, guard, xshift=-2mm] {\textit{isAdmin}} (inprog);

\draw[trans, red!60!black, bend left=35] (assigned) to node[below, guard, yshift=-2mm] {REASSIGN} node[above, guard, yshift=4mm] {\textit{isAdmin}} (open);

\end{tikzpicture}
\caption{Diagrama de la màquina d'estats XState per al cicle de vida de les incidències. Cada transició inclou l'esdeveniment (en majúscules) i el guard que controla el rol autoritzat (en cursiva).}
\label{fig:xstate}
\end{figure}


% ============================================================================
% DIAGRAMA 2: MIDDLEWARE — FLUX DE PETICIONS HTTP AMB JWT I RBAC
% ============================================================================

\begin{figure}[H]
\centering
\begin{tikzpicture}[
    node distance=1.4cm,
    block/.style={
        rectangle, rounded corners=4pt, draw=black, thick,
        minimum width=4.5cm, minimum height=0.9cm,
        font=\sffamily\small, align=center
    },
    decision/.style={
        diamond, draw=black, thick, aspect=2.5,
        minimum width=2cm, minimum height=0.8cm,
        font=\sffamily\small, align=center, inner sep=1pt
    },
    error/.style={
        rectangle, rounded corners=4pt, draw=red!70, fill=red!8,
        minimum width=2.8cm, minimum height=0.7cm,
        font=\sffamily\small, text=red!70!black
    },
    success/.style={
        rectangle, rounded corners=4pt, draw=green!60!black, fill=green!8,
        minimum width=3.5cm, minimum height=0.7cm,
        font=\sffamily\small, text=green!50!black
    },
    arr/.style={-Stealth, thick},
    label/.style={font=\sffamily\scriptsize}
]

% Flux principal
\node[block, fill=blue!8] (req) {Client HTTP Request};
\node[block, below=of req] (json) {express.json()};
\node[block, below=of json] (router) {Router (ruta + mètode)};
\node[decision, below=1.6cm of router] (protegida) {Ruta\\protegida?};

% Branca pública
\node[block, left=3cm of protegida, fill=green!8, yshift=-1cm] (ctrl_pub) {Controller};

% Branca autenticada
\node[block, below=1.8cm of protegida] (auth) {\texttt{authenticate}};
\node[decision, below=1.6cm of auth] (jwt_ok) {JWT\\vàlid?};
\node[error, left=3cm of jwt_ok] (e401) {401 Unauthorized};

\node[decision, below=1.6cm of jwt_ok] (rbac) {RBAC\\authorize?};
\node[error, left=3cm of rbac] (e403) {403 Forbidden};

\node[block, below=1.6cm of rbac, fill=green!8] (ctrl) {Controller};
\node[block, below=of ctrl, fill=yellow!10] (svc) {Service (lògica de negoci)};
\node[block, below=of svc, fill=gray!10] (db) {Prisma → PostgreSQL};
\node[block, below=of db, fill=blue!8] (res) {HTTP Response};

% Fletxes
\draw[arr] (req) -- (json);
\draw[arr] (json) -- (router);
\draw[arr] (router) -- (protegida);

\draw[arr] (protegida) -- node[left, label] {Sí} (auth);
\draw[arr] (protegida) -| node[above, label] {No} (ctrl_pub);
\draw[arr] (ctrl_pub) |- (svc);

\draw[arr] (auth) -- (jwt_ok);
\draw[arr] (jwt_ok) -- node[above, label] {No} (e401);
\draw[arr] (jwt_ok) -- node[left, label] {Sí} (rbac);
\draw[arr] (rbac) -- node[above, label] {No} (e403);
\draw[arr] (rbac) -- node[left, label] {Sí} (ctrl);
\draw[arr] (ctrl) -- (svc);
\draw[arr] (svc) -- (db);
\draw[arr] (db) -- (res);

\end{tikzpicture}
\caption{Flux complet d'una petició HTTP a través de la cadena de middlewares. Les rutes públiques (registre, login) eviten la capa d'autenticació, mentre que les rutes protegides passen per la verificació JWT i, opcionalment, pel control d'accés RBAC.}
\label{fig:middleware}
\end{figure}


% ============================================================================
% DIAGRAMA 3: FLUX D'AUTENTICACIÓ — REGISTRE, LOGIN, HASH I JWT
% ============================================================================

\begin{figure}[H]
\centering
\begin{tikzpicture}[
    node distance=1.2cm and 2.5cm,
    block/.style={
        rectangle, rounded corners=4pt, draw=black, thick,
        minimum width=3.5cm, minimum height=0.8cm,
        font=\sffamily\small, align=center
    },
    crypto/.style={
        block, fill=purple!10, draw=purple!50
    },
    db/.style={
        block, fill=gray!10, draw=gray!50
    },
    arr/.style={-Stealth, thick},
    label/.style={font=\sffamily\scriptsize},
    title/.style={font=\sffamily\bfseries\small}
]

% --- REGISTRE ---
\node[title] (t1) {REGISTRE};
\node[block, below=0.6cm of t1] (r1) {Client envia email,\\password, dades};
\node[crypto, below=of r1] (r2) {bcrypt.hash(password, 10)\\{\scriptsize genera salt aleatori}};
\node[db, below=of r2] (r3) {prisma.user.create(\{\\password: hash\})};
\node[block, below=of r3, fill=green!8] (r4) {Retorna usuari\\{\scriptsize (sense password)}};

\draw[arr] (r1) -- (r2);
\draw[arr] (r2) -- (r3);
\draw[arr] (r3) -- (r4);

% --- LOGIN ---
\node[title, right=5cm of t1] (t2) {LOGIN};
\node[block, below=0.6cm of t2] (l1) {Client envia\\email, password};
\node[db, below=of l1] (l2) {prisma.user.findUnique\\(\{email\})};
\node[crypto, below=of l2] (l3) {bcrypt.compare\\(password, hash)};
\node[crypto, below=of l3] (l4) {jwt.sign(\{userId, role\},\\secret, \{expiresIn: 24h\})};
\node[block, below=of l4, fill=green!8] (l5) {Retorna token\\+ dades usuari};

\draw[arr] (l1) -- (l2);
\draw[arr] (l2) -- (l3);
\draw[arr] (l3) -- (l4);
\draw[arr] (l4) -- (l5);

% --- VERIFICACIÓ ---
\node[title, right=5cm of t2] (t3) {VERIFICACIÓ};
\node[block, below=0.6cm of t3] (v1) {Petició amb header\\Authorization: Bearer \textless token\textgreater};
\node[crypto, below=of v1] (v2) {jwt.verify(token, secret)};
\node[block, below=of v2, fill=green!8] (v3) {Payload: \{userId, role\}\\injectat a req.user};
\node[block, below=of v3] (v4) {Continua al\\Controller};

\draw[arr] (v1) -- (v2);
\draw[arr] (v2) -- (v3);
\draw[arr] (v3) -- (v4);

\end{tikzpicture}
\caption{Flux criptogràfic dels tres processos d'autenticació. El registre genera un hash bcrypt amb salt aleatori, el login verifica el hash i genera un token JWT, i la verificació descodifica el token per injectar la identitat de l'usuari a cada petició.}
\label{fig:auth}
\end{figure}


% ============================================================================
% DIAGRAMA 4: ARQUITECTURA GENERAL DEL SISTEMA
% ============================================================================

\begin{figure}[H]
\centering
\begin{tikzpicture}[
    node distance=1.5cm and 2.5cm,
    component/.style={
        rectangle, rounded corners=6pt, draw=black, thick,
        minimum width=3.2cm, minimum height=1.8cm,
        font=\sffamily\small, align=center
    },
    arr/.style={-Stealth, thick},
    bidir/.style={Stealth-Stealth, thick},
    label/.style={font=\sffamily\scriptsize, text=gray!60!black}
]

% Frontend Web
\node[component, fill=blue!10, draw=blue!50] (web) {
    \textbf{Frontend Web}\\
    {\scriptsize React + Vite}\\
    {\scriptsize Tailwind CSS}\\
    {\scriptsize Recharts + Leaflet}
};

% Frontend Mòbil
\node[component, fill=cyan!10, draw=cyan!50, below=of web] (mobile) {
    \textbf{Frontend Mòbil}\\
    {\scriptsize (Sprint futur)}\\
    {\scriptsize React Native /}\\
    {\scriptsize Flutter}
};

% Backend
\node[component, fill=orange!10, draw=orange!50, right=3cm of web, yshift=-1.2cm] (backend) {
    \textbf{Backend API}\\
    {\scriptsize Express v5}\\
    {\scriptsize TypeScript}\\
    {\scriptsize XState + JWT}
};

% Supabase
\node[component, fill=green!10, draw=green!50, right=3cm of backend] (supabase) {
    \textbf{Supabase}\\
    {\scriptsize PostgreSQL}\\
    {\scriptsize PostGIS}\\
    {\scriptsize Auth + Storage}
};

% Fletxes
\draw[bidir, blue!60] (web.east) -- node[above, label] {REST API} node[below, label] {JSON + JWT} (web.east -| backend.west) -- (backend.west |- web.east);
\draw[bidir, cyan!60] (mobile.east) -- (mobile.east -| backend.west) -- (backend.west |- mobile.east);
\draw[bidir, green!60] (backend.east) -- node[above, label] {Prisma ORM} node[below, label] {SQL + PostGIS} (supabase.west);

% Etiqueta proxy
\node[label, above=0.3cm of web.north east, xshift=1cm] {Vite Proxy (:5173 → :3000)};

\end{tikzpicture}
\caption{Arquitectura general del sistema CityFix. El frontend web es comunica amb el backend via API REST, i el backend gestiona la persistència mitjançant Prisma ORM contra PostgreSQL allotjat a Supabase.}
\label{fig:arquitectura}
\end{figure}


% ============================================================================
% DIAGRAMA 5: SISTEMA D'INVITACIONS — REGISTRE PRIVILEGIAT
% ============================================================================

\begin{figure}[H]
\centering
\begin{tikzpicture}[
    node distance=1.3cm,
    block/.style={
        rectangle, rounded corners=4pt, draw=black, thick,
        minimum width=5cm, minimum height=0.8cm,
        font=\sffamily\small, align=center
    },
    decision/.style={
        diamond, draw=black, thick, aspect=2.5,
        minimum width=2cm, font=\sffamily\small,
        align=center, inner sep=1pt
    },
    error/.style={
        rectangle, rounded corners=4pt, draw=red!70, fill=red!8,
        minimum width=2.5cm, minimum height=0.7cm,
        font=\sffamily\small, text=red!70!black
    },
    success/.style={
        rectangle, rounded corners=4pt, draw=green!60!black, fill=green!8,
        minimum width=4cm, minimum height=0.7cm,
        font=\sffamily\small
    },
    arr/.style={-Stealth, thick},
    label/.style={font=\sffamily\scriptsize}
]

\node[block, fill=blue!8] (req) {POST /api/auth/register\\{\scriptsize \{email, password, role?, token?\}}};
\node[decision, below=1.5cm of req] (isPriv) {role és\\ADMIN o\\TECHNICAL?};

% Branca STUDENT
\node[block, left=3.5cm of isPriv, fill=green!8, yshift=-1cm] (force) {Força role = STUDENT\\{\scriptsize (ignora camp role)}};
\node[success, below=of force] (create_s) {prisma.user.create()};

% Branca privilegiada
\node[block, below=1.8cm of isPriv] (check) {Cerca invitació:\\{\scriptsize \{email, token, status: PENDING\}}};
\node[decision, below=1.5cm of check] (exists) {Invitació\\vàlida?};
\node[error, left=3.5cm of exists] (e403) {403 Forbidden};
\node[block, below=1.8cm of exists, fill=purple!8] (tx) {\$transaction: [\\user.create + invite.update]};
\node[success, below=of tx] (done) {Usuari creat + invitació USED};

\draw[arr] (req) -- (isPriv);
\draw[arr] (isPriv) -| node[above, label] {No} (force);
\draw[arr] (force) -- (create_s);
\draw[arr] (isPriv) -- node[left, label] {Sí} (check);
\draw[arr] (check) -- (exists);
\draw[arr] (exists) -- node[above, label] {No} (e403);
\draw[arr] (exists) -- node[left, label] {Sí} (tx);
\draw[arr] (tx) -- (done);

\end{tikzpicture}
\caption{Flux de registre amb el sistema d'invitacions. Els estudiants es registren lliurement (el backend força el rol), mentre que els rols privilegiats requereixen un token d'invitació vàlid generat per un administrador. La creació d'usuari i el consum de la invitació s'executen en transacció atòmica.}
\label{fig:invitacions}
\end{figure}


% ============================================================================
% DIAGRAMA 6: ARQUITECTURA BACKEND — PATRÓ PER CAPES
% ============================================================================

\begin{figure}[H]
\centering
\begin{tikzpicture}[
    layer/.style={
        rectangle, rounded corners=6pt, draw=#1!60, fill=#1!8, thick,
        minimum width=12cm, minimum height=1.3cm,
        font=\sffamily\small, align=center
    },
    arr/.style={-Stealth, thick, gray!60},
    note/.style={font=\sffamily\scriptsize, text=gray!50!black, align=left}
]

\node[layer=blue]   (client)  at (0, 7)   {\textbf{Client} (Navegador / App Mòbil)};
\node[layer=cyan]   (routes)  at (0, 5.2) {\textbf{Routes} — Defineix URL + mètode HTTP};
\node[layer=purple] (middle)  at (0, 3.4) {\textbf{Middlewares} — authenticate (JWT) + authorize (RBAC)};
\node[layer=orange] (ctrl)    at (0, 1.6) {\textbf{Controllers} — Valida entrada, retorna HTTP response};
\node[layer=yellow!80!orange] (svc) at (0, -0.2) {\textbf{Services} — Lògica de negoci (bcrypt, XState, queries)};
\node[layer=green!60!black]  (db)  at (0, -2)   {\textbf{Database} — Prisma Client → PostgreSQL (Supabase)};

\draw[arr] (client) -- (routes);
\draw[arr] (routes) -- (middle);
\draw[arr] (middle) -- (ctrl);
\draw[arr] (ctrl)   -- (svc);
\draw[arr] (svc)    -- (db);

% Anotacions laterals
\node[note, right=0.3cm of routes.east]  {auth.ts, reports.ts, users.ts,\\invites.ts, geo.ts, analytics.ts};
\node[note, right=0.3cm of middle.east]  {jwt.verify() → req.user\\authorize('ADMIN')};
\node[note, right=0.3cm of ctrl.east]    {Validació de body/params,\\gestió d'errors HTTP};
\node[note, right=0.3cm of svc.east]     {Transaccions, XState actors,\\agregacions, GeoJSON};
\node[note, right=0.3cm of db.east]      {Prisma v7 + PrismaPg adapter\\PostGIS extension};

\end{tikzpicture}
\caption{Arquitectura per capes del backend. Cada capa té una responsabilitat única i comunica exclusivament amb les capes adjacents, facilitant el manteniment i la testabilitat del sistema.}
\label{fig:capes}
\end{figure}


% ============================================================================
% DIAGRAMA 7: TRANSICIÓ D'ESTAT AMB XSTATE — FLUX DETALLAT
% ============================================================================

\begin{figure}[H]
\centering
\begin{tikzpicture}[
    node distance=1.2cm,
    block/.style={
        rectangle, rounded corners=4pt, draw=black, thick,
        minimum width=5.5cm, minimum height=0.8cm,
        font=\sffamily\small, align=center
    },
    decision/.style={
        diamond, draw=black, thick, aspect=2.2,
        minimum width=2.5cm, font=\sffamily\small,
        align=center, inner sep=1pt
    },
    error/.style={
        rectangle, rounded corners=4pt, draw=red!70, fill=red!8,
        minimum width=3.5cm, font=\sffamily\small, text=red!70!black
    },
    arr/.style={-Stealth, thick},
    label/.style={font=\sffamily\scriptsize}
]

\node[block, fill=blue!8] (req) {PATCH /reports/:id/transition\\{\scriptsize \{event: "ASSIGN", assignedToId?\}}};
\node[block, below=of req] (find) {prisma.report.findUnique(id)\\{\scriptsize obté estat actual de la BD}};
\node[block, below=of find, fill=purple!8] (actor) {createActor(stateMachine)\\{\scriptsize restaura estat amb resolveState()}};
\node[decision, below=1.5cm of actor] (can) {snapshot.can\\(\{type: event\})?};
\node[error, left=3.5cm of can] (err) {400: "Transició 'X' no\\permesa des de 'Y'\\amb rol 'Z'"};
\node[block, below=1.8cm of can, fill=green!8] (send) {actor.send(\{type: event\})\\{\scriptsize obté nou estat del snapshot}};
\node[block, below=of send] (update) {prisma.report.update\\(\{state: nouEstat\})};
\node[block, below=of update, fill=blue!8] (res) {200: Report actualitzat};

\draw[arr] (req) -- (find);
\draw[arr] (find) -- (actor);
\draw[arr] (actor) -- (can);
\draw[arr] (can) -- node[above, label] {No} (err);
\draw[arr] (can) -- node[left, label] {Sí} (send);
\draw[arr] (send) -- (update);
\draw[arr] (update) -- (res);

\end{tikzpicture}
\caption{Flux detallat d'una transició d'estat. El servei restaura l'estat actual des de la BD, valida la transició amb XState abans d'executar-la, i actualitza la BD només si la màquina ho autoritza.}
\label{fig:transicio}
\end{figure}


% ============================================================================
% DIAGRAMA 8: MODEL ENTITAT-RELACIÓ SIMPLIFICAT
% ============================================================================

\begin{figure}[H]
\centering
\begin{tikzpicture}[
    entity/.style={
        rectangle, rounded corners=4pt, draw=black, thick,
        minimum width=3.5cm, font=\sffamily\small, align=left,
        inner sep=8pt
    },
    arr/.style={-Stealth, thick},
    rel/.style={font=\sffamily\scriptsize, text=gray!60!black}
]

% Entitats
\node[entity, fill=blue!8] (user) {
    \textbf{User}\\
    \rule{2.5cm}{0.4pt}\\
    {\scriptsize user\_id (PK)}\\
    {\scriptsize email, name, surname}\\
    {\scriptsize nickname, password}\\
    {\scriptsize role: enum}\\
    {\scriptsize active: boolean}\\
    {\scriptsize points: int}\\
    {\scriptsize inviteId? (FK)}
};

\node[entity, fill=orange!8, right=4cm of user] (report) {
    \textbf{Report}\\
    \rule{2.5cm}{0.4pt}\\
    {\scriptsize report\_id (PK)}\\
    {\scriptsize title, description}\\
    {\scriptsize state: enum}\\
    {\scriptsize priority: enum}\\
    {\scriptsize category: enum}\\
    {\scriptsize latitude, longitude}\\
    {\scriptsize createdById (FK)}\\
    {\scriptsize assignedToId? (FK)}\\
    {\scriptsize resolvedAt?}
};

\node[entity, fill=green!8, below=3cm of user] (invite) {
    \textbf{Invite}\\
    \rule{2.5cm}{0.4pt}\\
    {\scriptsize id (PK)}\\
    {\scriptsize email, role}\\
    {\scriptsize token (unique)}\\
    {\scriptsize status: enum}\\
    {\scriptsize createdAt}
};

\node[entity, fill=purple!8, below=2.5cm of report, xshift=-0.5cm] (image) {
    \textbf{Image}\\
    \rule{2cm}{0.4pt}\\
    {\scriptsize id (PK)}\\
    {\scriptsize url}\\
    {\scriptsize reportId (FK)}
};

\node[entity, fill=yellow!10, below=2.5cm of report, xshift=3.5cm] (comment) {
    \textbf{Comment}\\
    \rule{2cm}{0.4pt}\\
    {\scriptsize id (PK)}\\
    {\scriptsize content}\\
    {\scriptsize userId (FK)}\\
    {\scriptsize reportId (FK)}
};

% Relacions
\draw[arr, blue!60] (user.east) -- node[above, rel] {crea (1:N)} (report.west);
\draw[arr, cyan!60] ([yshift=-5mm]user.east) -- node[below, rel] {assignat (1:N)} ([yshift=-5mm]report.west);

\draw[arr, green!60] (invite.north) -- node[left, rel] {origina (1:1)} (user.south);

\draw[arr, purple!60] (report.south) -- node[left, rel] {té (1:N)} (image.north);
\draw[arr, yellow!60!black] ([xshift=5mm]report.south) -- node[right, rel] {té (1:N)} (comment.north);

\draw[arr, gray] (user.south east) |- node[below, rel, near end] {escriu (1:N)} (comment.west);

\end{tikzpicture}
\caption{Model Entitat-Relació simplificat del sistema CityFix. Un usuari pot crear múltiples incidències i ser assignat com a tècnic. Les invitacions vinculen cada accés privilegiat al seu origen.}
\label{fig:er}
\end{figure}


% ============================================================================
% DIAGRAMA 9: FLUX DE REVOCACIÓ AMB REGLES ANTI-DEADLOCK
% ============================================================================

\begin{figure}[H]
\centering
\begin{tikzpicture}[
    node distance=1.3cm,
    block/.style={
        rectangle, rounded corners=4pt, draw=black, thick,
        minimum width=4.5cm, minimum height=0.8cm,
        font=\sffamily\small, align=center
    },
    decision/.style={
        diamond, draw=black, thick, aspect=2.2,
        minimum width=2cm, font=\sffamily\small,
        align=center, inner sep=1pt
    },
    error/.style={
        rectangle, rounded corners=4pt, draw=red!70, fill=red!8,
        minimum width=3cm, minimum height=0.7cm,
        font=\sffamily\small, text=red!70!black
    },
    arr/.style={-Stealth, thick},
    label/.style={font=\sffamily\scriptsize}
]

\node[block, fill=blue!8] (req) {PATCH /users/:id/revoke};
\node[decision, below=1.5cm of req] (root) {És el\\root admin?};
\node[error, right=3.5cm of root] (e1) {403: Root admin\\intocable};
\node[decision, below=1.8cm of root] (last) {Últim admin\\actiu?};
\node[error, right=3.5cm of last] (e2) {403: No es pot\\revocar l'últim};
\node[block, below=1.8cm of last, fill=purple!8] (tx) {\$transaction: [\\user.active=false,\\invite.status=REVOKED]};
\node[block, below=of tx, fill=green!8] (done) {Usuari revocat\\{\scriptsize (no pot fer login)}};

\draw[arr] (req) -- (root);
\draw[arr] (root) -- node[above, label] {Sí} (e1);
\draw[arr] (root) -- node[left, label] {No} (last);
\draw[arr] (last) -- node[above, label] {Sí} (e2);
\draw[arr] (last) -- node[left, label] {No} (tx);
\draw[arr] (tx) -- (done);

\end{tikzpicture}
\caption{Flux de revocació d'un usuari privilegiat amb les dues regles anti-deadlock. La regla 1 protegeix el root admin i la regla 2 impedeix eliminar l'últim administrador actiu.}
\label{fig:revocacio}
\end{figure}
